\textbf{Phần 1. Trắc nghiệm lựa chọn}
\setcounterpageref{ex}{0}

\begin{ex} %[2H3V15]
Trên đĩa cân bên trái có một quả cân $5\text{ kg}$, đĩa bên phải có $1$ gói bột giặt và $2$ quả cân $1\text{ kg}$. Biết cân thăng bằng, gói bột giặt nặng bao nhiêu ki-lô-gam?
\choice
{\True $3\text{ kg}$}
{$2\text{ kg}$}
{$4\text{ kg}$}
{$5\text{ kg}$}
\loigiai{
Vì cân thăng bằng nên: Khối lượng gói bột giặt $+ 1 + 1 = 5$. Khối lượng gói bột giặt là: $5 - 2 = 3\text{ (kg)}$.
\textbf{Đáp án của câu hỏi là A}
}
\end{ex}

\begin{ex} %[2H3V15]
Bao gạo nặng $35\text{ kg}$, bao ngô nhẹ hơn bao gạo $8\text{ kg}$. Hỏi cả hai bao nặng bao nhiêu ki-lô-gam?
\choice
{$54\text{ kg}$}
{$60\text{ kg}$}
{\True $62\text{ kg}$}
{$58\text{ kg}$}
\loigiai{
Bao ngô nặng: $35 - 8 = 27\text{ (kg)}$.\nCả hai bao nặng: $35 + 27 = 62\text{ (kg)}$.
\textbf{Đáp án của câu hỏi là C}
}
\end{ex}

\begin{ex} %[2H3V15]
Tính giá trị của biểu thức: $48\text{ kg} - 19\text{ kg} + 25\text{ kg}$.
\choice
{$52\text{ kg}$}
{\True $54\text{ kg}$}
{$56\text{ kg}$}
{$50\text{ kg}$}
\loigiai{
$48 - 19 = 29\text{ kg}$; $29 + 25 = 54\text{ kg}$.
\textbf{Đáp án của câu hỏi là B}
}
\end{ex}

\begin{ex} %[2H3V15]
Con lợn đen nặng $65\text{ kg}$, con lợn trắng nhẹ hơn con lợn đen $18\text{ kg}$. Con lợn trắng nặng bao nhiêu ki-lô-gam?
\choice
{$49\text{ kg}$}
{$53\text{ kg}$}
{$45\text{ kg}$}
{\True $47\text{ kg}$}
\loigiai{
Con lợn trắng nặng: $65 - 18 = 47\text{ (kg)}$.
\textbf{Đáp án của câu hỏi là D}
}
\end{ex}

\begin{ex} %[2H3V15]
Một thùng chứa đầy nước nặng $22\text{ kg}$. Sau khi người ta dùng đi một nửa số nước thì thùng và số nước còn lại nặng $12\text{ kg}$. Hỏi vỏ thùng không có nước nặng bao nhiêu ki-lô-gam?
\choice
{\True $2\text{ kg}$}
{$3\text{ kg}$}
{$1\text{ kg}$}
{$4\text{ kg}$}
\loigiai{
Nửa số nước nặng: $22 - 12 = 10\text{ (kg)}$.\nCả thùng nước chứa số nước nặng: $10 \times 2 = 20\text{ (kg)}$.\nVỏ thùng nặng: $22 - 20 = 2\text{ (kg)}$.
\textbf{Đáp án của câu hỏi là A}
}
\end{ex}

\begin{ex} %[2H3V15]
Điền dấu thích hợp: $36\text{ kg} + 28\text{ kg} \dots 90\text{ kg} - 25\text{ kg}$.
\choice
{$>$}
{$=$}
{\True $<$}
{Không so sánh được}
\loigiai{
Vế trái: $36 + 28 = 64\text{ kg}$. Vế phải: $90 - 25 = 65\text{ kg}$. Vì $64 < 65$ nên điền dấu $<$. 
\textbf{Đáp án của câu hỏi là C}
}
\end{ex}

\begin{ex} %[2H3V15]
Có ba bao thóc: bao thứ nhất nặng $24\text{ kg}$, bao thứ hai nặng $26\text{ kg}$, bao thứ ba nặng $30\text{ kg}$. Cả ba bao thóc nặng tất cả bao nhiêu ki-lô-gam?
\choice
{$78\text{ kg}$}
{\True $80\text{ kg}$}
{$82\text{ kg}$}
{$75\text{ kg}$}
\loigiai{
$(24 + 26) + 30 = 50 + 30 = 80\text{ (kg)}$.
\textbf{Đáp án của câu hỏi là B}
}
\end{ex}

\begin{ex} %[2H3V15]
Mẹ mua một con gà nặng $2\text{ kg}$ và một con vịt nặng $3\text{ kg}$. Bác An mua một con ngỗng nặng bằng cả con gà và con vịt cộng thêm $2\text{ kg}$. Con ngỗng nặng bao nhiêu?
\choice
{$6\text{ kg}$}
{$8\text{ kg}$}
{$5\text{ kg}$}
{\True $7\text{ kg}$}
\loigiai{
Con ngỗng nặng: $2 + 3 + 2 = 7\text{ (kg)}$.
\textbf{Đáp án của câu hỏi là D}
}
\end{ex}

\begin{ex} %[2H3V15]
Một xe chở $45\text{ kg}$ đường và $38\text{ kg}$ bột mì. Khối lượng hàng hóa xe chở là
\choice
{\True $83\text{ kg}$}
{$73\text{ kg}$}
{$81\text{ kg}$}
{$85\text{ kg}$}
\loigiai{
$45 + 38 = 83\text{ (kg)}$.
\textbf{Đáp án của câu hỏi là A}
}
\end{ex}

\begin{ex} %[2H3V15]
Con bò nặng hơn con bê $45\text{ kg}$. Biết con bò nặng $90\text{ kg}$, con bê nặng bao nhiêu ki-lô-gam?
\choice
{$55\text{ kg}$}
{$40\text{ kg}$}
{\True $45\text{ kg}$}
{$50\text{ kg}$}
\loigiai{
Con bê nặng: $90 - 45 = 45\text{ (kg)}$.
\textbf{Đáp án của câu hỏi là C}
}
\end{ex}

\begin{ex} %[2H3V15]
Túi muối nặng $2\text{ kg}$, túi đường nặng gấp đôi túi muối. Hỏi cả hai túi nặng bao nhiêu ki-lô-gam?
\choice
{$4\text{ kg}$}
{\True $6\text{ kg}$}
{$5\text{ kg}$}
{$7\text{ kg}$}
\loigiai{
Túi đường nặng: $2 + 2 = 4\text{ (kg)}$. Cả hai túi nặng: $2 + 4 = 6\text{ (kg)}$.
\textbf{Đáp án của câu hỏi là B}
}
\end{ex}

\begin{ex} %[2H3V15]
Tìm $x$ biết: $x + 27\text{ kg} = 65\text{ kg}$.
\choice
{$37\text{ kg}$}
{$42\text{ kg}$}
{$48\text{ kg}$}
{\True $38\text{ kg}$}
\loigiai{
$x = 65 - 27 = 38\text{ (kg)}$.
\textbf{Đáp án của câu hỏi là D}
}
\end{ex}

\textbf{Phần 2. Trắc nghiệm đúng sai}
\setcounterpageref{ex}{0}

\begin{ex} %[2H3V15]
Xét tính đúng sai của các bài toán về ki-lô-gam:
\choiceTF
{\True Ki-lô-gam là đơn vị dùng để đo khối lượng, viết tắt là kg.}
{Quả cân 5 kg nhẹ hơn tổng của hai quả cân 2 kg và một quả cân 1 kg.}
{\True Một túi gạo 10 kg san đều vào hai túi thì mỗi túi có 5 kg.}
{\True Khối lượng của hai vật bằng nhau thì đặt lên hai đĩa của cân đĩa, cân sẽ thăng bằng.}
\loigiai{
\begin{itemize}
    \item Ý a) Đúng: Quy ước chuẩn theo sách giáo khoa.
    \item Ý b) Sai: Hai quả 2 kg và một quả 1 kg là $2 + 2 + 1 = 5\text{ kg}$, bằng nhau.
    \item Ý c) Đúng: $10 : 2 = 5\text{ kg}$. Đúng.
    \item Ý d) Đúng: Nguyên lý cân đĩa thăng bằng.
\end{itemize}
\begin{itemize}
    \item a) \textbf{Đúng}.
    \item b) \textbf{Sai}.
    \item c) \textbf{Đúng}.
    \item d) \textbf{Đúng}.
\end{itemize}
}
\end{ex}

\begin{ex} %[2H3V15]
Bác Nam có hai bao thóc, bao thứ nhất nặng 42 kg, bao thứ hai nặng 38 kg:
\choiceTF
{\True Bao thứ nhất nặng hơn bao thứ hai 4 kg.}
{\True Cả hai bao nặng tất cả 80 kg.}
{Chuyển 4 kg từ bao thứ nhất sang bao thứ hai thì hai bao nặng bằng nhau.}
{Bao thứ hai nặng gấp đôi bao thứ nhất.}
\loigiai{
\begin{itemize}
    \item Ý a) Đúng: $42 - 38 = 4\text{ kg}$.
    \item Ý b) Đúng: $42 + 38 = 80\text{ kg}$.
    \item Ý c) Sai: Chỉ cần chuyển $4 : 2 = 2\text{ kg}$ thì hai bao đều nặng 40 kg.
    \item Ý d) Sai: 38 không gấp đôi 42. Khẳng định sai.
\end{itemize}
\begin{itemize}
    \item a) \textbf{Đúng}.
    \item b) \textbf{Đúng}.
    \item c) \textbf{Sai}.
    \item d) \textbf{Sai}.
\end{itemize}
}
\end{ex}

\textbf{Phần 3. Trả lời ngắn}
\setcounterpageref{ex}{0}

\begin{ex} %[2H3V15]
Mẹ mua về 24 kg gạo. Sau khi dùng một tuần thì còn lại 17 kg. Hỏi gia đình đã ăn hết bao nhiêu ki-lô-gam gạo?
\shortans{7}
\loigiai{
\textbf{Đáp số:} $7$\\
$24 - 17 = 7\text{ kg}$.
}
\end{ex}

\begin{ex} %[2H3V15]
Có hai quả dưa hấu, quả thứ nhất nặng 4 kg, quả thứ hai nặng hơn quả thứ nhất 2 kg. Hỏi cả hai quả dưa hấu nặng bao nhiêu ki-lô-gam?
\shortans{10}
\loigiai{
\textbf{Đáp số:} $10$\\
Quả 2: $4 + 2 = 6\text{ kg}$. Tổng: $4 + 6 = 10\text{ kg}$.
}
\end{ex}

\begin{ex} %[2H3V15]
Con bê nặng 35 kg, con lợn nặng 28 kg. Hỏi con bê nặng hơn con lợn bao nhiêu ki-lô-gam?
\shortans{7}
\loigiai{
\textbf{Đáp số:} $7$\\
$35 - 28 = 7\text{ kg}$.
}
\end{ex}

\begin{ex} %[2H3V15]
Một cửa hàng có 50 kg đường, đã bán đi 26 kg. Số đường còn lại chia đều vào 3 túi. Hỏi mỗi túi có bao nhiêu ki-lô-gam đường?
\shortans{8}
\loigiai{
\textbf{Đáp số:} $8$\\
Còn lại: $50 - 26 = 24\text{ kg}$. Mỗi túi: $24 : 3 = 8\text{ kg}$.
}
\end{ex}

\begin{ex} %[2H3V15]
Để cân một con gà, người ta đặt lên một đĩa con gà và quả cân 1 kg, đĩa kia đặt quả cân 4 kg thì cân thăng bằng. Hỏi con gà nặng bao nhiêu ki-lô-gam?
\shortans{3}
\loigiai{
\textbf{Đáp số:} $3$\\
$4 - 1 = 3\text{ kg}$.
}
\end{ex}

\begin{ex} %[2H3V15]
Tổng khối lượng của một bao gạo và một bao ngô là 65 kg, biết bao gạo nặng 37 kg. Hỏi bao ngô nặng bao nhiêu ki-lô-gam?
\shortans{28}
\loigiai{
\textbf{Đáp số:} $28$\\
$65 - 37 = 28\text{ kg}$.
}
\end{ex}
